\scanpage{13}
où \deleted{$\unsure{h_1}$ \unsure{est le}}
\[
(15)\qquad h_1=\operatorname{rg}_{\mathbb Z}H_1(G_X,\mathbb Z).
\]
\uline{Démonstration} Standart (!) -- la formule (13) résulte d'une suite exacte \inserted{standart} de cohomologie reliant $\underline O_{\widehat X}$, $\underline O_{\widetilde{\widehat X}}$ et les singularités de $\widehat X$, d'ailleurs la formule d'EP pour le graphe $G_X$ nous donne
\[
\chi(G_X)=\underbrace{h_0}_{1}\deleted{\unsure{(G_X)}}-h_1=\operatorname{card}S-\operatorname{card}A
\]
i.e.
\[
1-\operatorname{card}S+\operatorname{card}A=h_1\qquad\text{O.K.}
\]
\uline{Corollaire} Le nb $g(\widehat X)$ ne dépend que de la maquette de $X$, par la formule (14)
\[
g(\widehat X)=\underbrace{\sum_{\alpha\in S}g(\alpha)+h_1}_{\overset{\mathrm{dfn}}{=}\ \text{genre de }G}
\]
-- on l'appelle le \uline{genre} de la maquette.

On appelle \uline{type} de la maquette le couple $(\bar g,\nu)$, où $g$ est son genre et $\nu=\operatorname{card}I$ est le \uline{poids marqué}. On appelle \uline{poids marqué} d'un sommet l'entier naturel
\mspage{12}
